%==============================================================
% PARTE 2: Resolución de un Laberinto
%==============================================================
\chapter{Parte 2: Resolución de un Laberinto}

En esta parte se utiliza el Turtlebot 3 equipado con un sensor LiDAR para resolver
un laberinto de manera autónoma. Se emplea el paquete \texttt{maze\_pkg}, que contiene
el nodo \texttt{res\_maze.py} basado en un algoritmo de seguimiento de pared izquierda
(\textit{left-wall following}).

%--------------------------------------------------------------
\section{Preguntas teóricas previas}
%--------------------------------------------------------------

\subsection*{Pregunta 12: Topic y tipo de mensaje del LiDAR}

El topic asociado al LiDAR del Turtlebot 3 es \texttt{/scan}. \\ El tipo de mensaje
es \texttt{sensor\_msgs/msg/LaserScan}, que contiene los campos principales:
\begin{itemize}
  \item \texttt{angle\_min}, \texttt{angle\_max}, \texttt{angle\_increment}: rango y resolución angular.
  \item \texttt{range\_min}, \texttt{range\_max}: límites de distancia medibles.
  \item \texttt{ranges[]}: array con las distancias medidas en cada ángulo.
\end{itemize}

Para verificarlo se puede ejecutar:
\begin{lstlisting}[style=bash]
ros2 topic list
ros2 topic info /scan
ros2 topic echo /scan --once
\end{lstlisting}

\subsection*{Pregunta 13: Rangos de distancia, ángulo de escaneo y origen de referencia}

Los valores del LiDAR del Turtlebot 3 Waffle y Burger son:

\begin{table}[H]
\centering
\begin{tabular}{lcc}
\toprule
\textbf{Parámetro}         & \textbf{Waffle}   & \textbf{Burger}   \\
\midrule
Rango mínimo               & 0.12 m            & 0.12 m            \\
Rango máximo               & 3.5 m             & 3.5 m             \\
Ángulo de escaneo          & 360°              & 360°              \\
Resolución angular         & 1°                & 1°                \\
\bottomrule
\end{tabular}
\caption{Características del LiDAR del Turtlebot 3}
\end{table}

El origen de referencia del LiDAR es el propio sensor, montado en la parte
superior del robot. El índice 0 del array \texttt{ranges[]} corresponde a la
dirección de avance del robot (frente), y los índices aumentan en sentido
antihorario: índice 90 $\approx$ izquierda, 180 $\approx$ atrás, 270 $\approx$ derecha.

% ---- CAPTURA RECOMENDADA ----
% Añadir aquí una captura de: ros2 topic echo /scan --once
% mostrando los valores de range_min, range_max, angle_min, angle_max
\begin{figure}[H]
  \centering
  % \includegraphics[width=0.8\textwidth]{parts/parte2/images/scan_info.png}
  \fbox{\parbox{0.8\textwidth}{\centering\vspace{1cm}
    [CAPTURA: salida de \texttt{ros2 topic echo /scan --once} mostrando
    range\_min, range\_max, angle\_min, angle\_max]
  \vspace{1cm}}}
  \caption{Información del topic \texttt{/scan} del Turtlebot 3}
\end{figure}

%--------------------------------------------------------------
\section{Ejercicio 1: Resolución del Laberinto 1 (Waffle)}
%--------------------------------------------------------------

Se ha desarrollado el nodo \texttt{res\_maze.py} dentro del paquete \texttt{maze\_pkg}.
El algoritmo implementado es \textit{left-wall following} con una máquina de estados:

\begin{itemize}
  \item \textbf{FIND\_WALL}: el robot avanza en línea recta hasta detectar una pared al
        frente ($d_{\text{front}} < 0.45$ m). Este estado inicial es necesario para que
        el robot tenga una pared de referencia antes de comenzar el seguimiento.
  \item \textbf{TURN\_RIGHT}: el robot gira sobre sí mismo hacia la derecha hasta que el
        frente queda libre. Para evitar oscilaciones, existe un contador mínimo de
        iteraciones (\texttt{MIN\_TURN\_TICKS = 15}) antes de poder salir del estado.
  \item \textbf{FORWARD}: el robot avanza siguiendo la pared izquierda mediante un
        controlador proporcional. El error se define como
        $e = d_{\text{left}} - d_{\text{desired}}$, de forma que si la pared está
        lejos el robot gira a la izquierda y si está cerca gira a la derecha.
\end{itemize}

Los parámetros principales del nodo son:

\begin{table}[H]
\centering
\begin{tabular}{lll}
\toprule
\textbf{Parámetro}    & \textbf{Valor} & \textbf{Descripción} \\
\midrule
\texttt{LINEAR\_SPEED}   & 0.12 m/s & Velocidad lineal de avance \\
\texttt{ANGULAR\_SPEED}  & 0.40 rad/s & Velocidad angular de giro \\
\texttt{WALL\_DIST}      & 0.35 m & Distancia deseada a la pared izquierda \\
\texttt{FRONT\_THRESH}   & 0.45 m & Umbral de obstáculo al frente \\
\texttt{KP}              & 0.8    & Ganancia del controlador proporcional \\
\texttt{MIN\_TURN\_TICKS}& 15     & Iteraciones mínimas en estado de giro \\
\bottomrule
\end{tabular}
\caption{Parámetros del nodo \texttt{res\_maze.py}}
\end{table}

Para lanzar el ejercicio:
\begin{lstlisting}[style=bash]
# Terminal 1 - Simulacion
export TURTLEBOT3_MODEL=waffle
ros2 launch maze_pkg maze_1.launch.py

# Terminal 2 - Solver
source /workspace/ros2_ws/install/setup.bash
ros2 run maze_pkg res_maze
\end{lstlisting}

% ---- CAPTURAS RECOMENDADAS ----
\begin{figure}[H]
  \centering
  % \includegraphics[width=0.8\textwidth]{parts/parte2/images/maze1_waffle_inicio.png}
  \fbox{\parbox{0.8\textwidth}{\centering\vspace{1cm}
    [CAPTURA: Gazebo con el Turtlebot 3 Waffle en el laberinto 1 (estado inicial)]
  \vspace{1cm}}}
  \caption{Turtlebot 3 Waffle en el laberinto 1 — estado inicial}
\end{figure}

\begin{figure}[H]
  \centering
  % \includegraphics[width=0.8\textwidth]{parts/parte2/images/maze1_waffle_resuelto.png}
  \fbox{\parbox{0.8\textwidth}{\centering\vspace{1cm}
    [CAPTURA: Gazebo con el Turtlebot 3 Waffle habiendo salido del laberinto 1]
  \vspace{1cm}}}
  \caption{Turtlebot 3 Waffle resolviendo el laberinto 1}
\end{figure}

El enlace al vídeo de la ejecución se puede encontrar en: \url{ENLACE_VIDEO_MAZE1_WAFFLE}

%--------------------------------------------------------------
\section{Ejercicio 2: Laberinto 2 y análisis de limitaciones}
%--------------------------------------------------------------

Se ejecuta el mismo nodo \texttt{res\_maze.py} con el segundo mapa:

\begin{lstlisting}[style=bash]
export TURTLEBOT3_MODEL=waffle
ros2 launch maze_pkg maze_2.launch.py
\end{lstlisting}

% ---- CAPTURA RECOMENDADA ----
\begin{figure}[H]
  \centering
  % \includegraphics[width=0.8\textwidth]{parts/parte2/images/maze2_waffle.png}
  \fbox{\parbox{0.8\textwidth}{\centering\vspace{1cm}
    [CAPTURA: Gazebo con el Turtlebot 3 Waffle en el laberinto 2]
  \vspace{1cm}}}
  \caption{Turtlebot 3 Waffle en el laberinto 2}
\end{figure}

\subsection*{Pregunta 14: Problemática observada en el laberinto 2}

El laberinto 2 presenta una topología diferente al primero. La principal problemática
observada es la presencia de \textbf{bucles cerrados} y pasillos simétricos que pueden
llevar al robot a girar indefinidamente por la misma ruta sin encontrar la salida.
El algoritmo de seguimiento de pared izquierda es completo para laberintos simplemente
conexos (sin islas interiores), pero falla cuando el laberinto contiene paredes
desconectadas del perímetro, ya que el robot puede quedar atrapado orbitando
alrededor de una isla interior.

\subsection*{Pregunta 15: ¿Es el robot capaz de resolver este laberinto?}

% TODO: completar tras ejecutar el experimento
El robot \textbf{[sí/no] es capaz} de resolver el laberinto 2 con el algoritmo
de seguimiento de pared izquierda. [Completar con la observación real tras la ejecución.]

En caso de no poder resolverlo, la información adicional que necesitaría el robot
para lograrlo sería:
\begin{itemize}
  \item Un \textbf{mapa previo} del entorno (navegación basada en mapa, SLAM).
  \item \textbf{Localización global} para detectar si ha visitado una posición anteriormente
        y así evitar repetir rutas (algoritmo de Pledge o similares).
  \item Capacidad de \textbf{marcar el camino recorrido} para implementar estrategias
        como \textit{Trémaux's algorithm}.
\end{itemize}

El enlace al vídeo de la ejecución se puede encontrar en: \url{ENLACE_VIDEO_MAZE2_WAFFLE}

%--------------------------------------------------------------
\section{Ejercicio 3: Repetición con Turtlebot 3 Burger}
%--------------------------------------------------------------

Se repiten los ejercicios anteriores usando el modelo \texttt{Burger}:

\begin{lstlisting}[style=bash]
export TURTLEBOT3_MODEL=burger
ros2 launch maze_pkg maze_1.launch.py
# Terminal 2:
ros2 run maze_pkg res_maze
\end{lstlisting}

% ---- CAPTURAS RECOMENDADAS ----
\begin{figure}[H]
  \centering
  % \includegraphics[width=0.8\textwidth]{parts/parte2/images/maze1_burger.png}
  \fbox{\parbox{0.8\textwidth}{\centering\vspace{1cm}
    [CAPTURA: Gazebo con Turtlebot 3 Burger en el laberinto 1]
  \vspace{1cm}}}
  \caption{Turtlebot 3 Burger en el laberinto 1}
\end{figure}

\subsection*{Pregunta 16: Diferencias observadas entre Waffle y Burger en el laberinto}

Las diferencias observadas entre ambos modelos al resolver el laberinto son las siguientes:

\begin{itemize}
  \item \textbf{Dimensiones físicas}: el Burger es más pequeño y estrecho que el Waffle,
        lo que le permite navegar con mayor holgura por pasillos estrechos y cometer
        menos colisiones laterales.
  \item \textbf{LiDAR}: ambos modelos montan el mismo sensor LiDAR con 360° de cobertura,
        por lo que la información de percepción es equivalente.
  \item \textbf{Dinámica de giro}: el Burger, al tener menor separación entre ruedas
        (\textit{wheel separation}), realiza giros más cerrados. Esto puede ser ventajoso
        en espacios reducidos pero también más sensible a errores en la odometría.
  \item \textbf{Velocidad y estabilidad}: el Waffle, al ser más pesado y ancho, es más
        estable a altas velocidades. El Burger tiende a oscilar más con los mismos
        parámetros del controlador.
  \item \textbf{Comportamiento en el laberinto}: % TODO completar con observación real
        [Completar tras la ejecución — indicar si el Burger resuelve el laberinto
        con mayor o menor facilidad, si choca con paredes, si los parámetros necesitan
        ajuste, etc.]
\end{itemize}

El enlace al vídeo de la ejecución con Burger se puede encontrar en: \url{ENLACE_VIDEO_BURGER}

%--------------------------------------------------------------
\section{Ejercicio 4: Laberinto propio}
%--------------------------------------------------------------

Se ha generado un laberinto personalizado usando el generador online
\href{https://mazegenerator.net/}{mazegenerator.net} y se ha implementado como
modelo SDF en Gazebo.

% ---- CAPTURA RECOMENDADA ----
\begin{figure}[H]
  \centering
  % \includegraphics[width=0.6\textwidth]{parts/parte2/images/maze_custom_generador.png}
  \fbox{\parbox{0.8\textwidth}{\centering\vspace{1cm}
    [CAPTURA: imagen del laberinto generado en mazegenerator.net]
  \vspace{1cm}}}
  \caption{Laberinto personalizado generado con mazegenerator.net}
\end{figure}

Para lanzar el laberinto personalizado:
\begin{lstlisting}[style=bash]
# Con Waffle
export TURTLEBOT3_MODEL=waffle
ros2 launch maze_pkg maze_3.launch.py

# Con Burger
export TURTLEBOT3_MODEL=burger
ros2 launch maze_pkg maze_3.launch.py
\end{lstlisting}

\begin{figure}[H]
  \centering
  % \includegraphics[width=0.8\textwidth]{parts/parte2/images/maze3_waffle.png}
  \fbox{\parbox{0.8\textwidth}{\centering\vspace{1cm}
    [CAPTURA: Gazebo con el laberinto personalizado y el Waffle]
  \vspace{1cm}}}
  \caption{Turtlebot 3 Waffle resolviendo el laberinto personalizado}
\end{figure}

\begin{figure}[H]
  \centering
  % \includegraphics[width=0.8\textwidth]{parts/parte2/images/maze3_burger.png}
  \fbox{\parbox{0.8\textwidth}{\centering\vspace{1cm}
    [CAPTURA: Gazebo con el laberinto personalizado y el Burger]
  \vspace{1cm}}}
  \caption{Turtlebot 3 Burger resolviendo el laberinto personalizado}
\end{figure}

El enlace al vídeo de la ejecución con el laberinto personalizado:
\url{ENLACE_VIDEO_MAZE3}
